% Standalone problem document, generated from the adjacent README.md.
% Compile with XeLaTeX. Mathematical content is maintained in Markdown.
\documentclass[11pt,a4paper]{article}
\usepackage[fontsize=10.5pt]{scrextend}
\usepackage[margin=22mm,headheight=15pt,headsep=9mm,footskip=11mm]{geometry}
\usepackage{fontspec}
\setmainfont{texgyrepagella}[Extension=.otf,UprightFont=*-regular,BoldFont=*-bold,ItalicFont=*-italic,BoldItalicFont=*-bolditalic]
\setsansfont{texgyreheros}[Extension=.otf,UprightFont=*-regular,BoldFont=*-bold,ItalicFont=*-italic,BoldItalicFont=*-bolditalic]
\setmonofont{texgyrecursor}[Extension=.otf,UprightFont=*-regular,BoldFont=*-bold,ItalicFont=*-italic,BoldItalicFont=*-bolditalic,Scale=MatchLowercase]
\usepackage{amsmath,amssymb}
\usepackage{unicode-math}
\setmathfont{texgyrepagella-math.otf}
\usepackage{xcolor,microtype,enumitem,needspace,fancyhdr}
\usepackage{bookmark}
\definecolor{ink}{HTML}{17344B}
\definecolor{accent}{HTML}{197D80}
\definecolor{muted}{HTML}{596773}
\hypersetup{colorlinks=true,linkcolor=accent,urlcolor=accent,pdftitle={TR-18 - Near-optimal
type-2 bound for Gaussian sums of symmetric
tensors},pdfauthor={Open Problems in Numerical Linear Algebra}}
\urlstyle{same}
\setlength{\parindent}{0pt}
\setlength{\parskip}{5pt plus 1pt minus 1pt}
\linespread{1.04}
\setlength{\emergencystretch}{3em}
\widowpenalty=10000
\clubpenalty=10000
\displaywidowpenalty=10000
\allowdisplaybreaks[1]
\setlist{nosep,leftmargin=1.5em}
\providecommand{\tightlist}{\setlength{\itemsep}{0pt}\setlength{\parskip}{0pt}}
\providecommand{\pandocbounded}[1]{#1}
\pagestyle{fancy}
\fancyhf{}
\fancyhead[L]{\sffamily\footnotesize\color{muted}OPEN PROBLEMS IN NUMERICAL LINEAR ALGEBRA}
\fancyhead[R]{\sffamily\bfseries\footnotesize\color{accent}TR-18}
\fancyfoot[L]{\parbox[b]{0.9\textwidth}{\sffamily\fontsize{8}{10}\selectfont\color{muted}Editorial ratings; a bounded search cannot certify that a problem remains unsolved.\\Literature check: 2026-09-10}}
\fancyfoot[R]{\sffamily\footnotesize\color{muted}\thepage}
\renewcommand{\headrulewidth}{0.3pt}
\renewcommand{\headrule}{\hbox to\headwidth{\color{accent}\leaders\hrule height \headrulewidth\hfill}}
\makeatletter
\renewcommand{\section}{\Needspace{5\baselineskip}\@startsection{section}{1}{0pt}{12pt plus 2pt minus 2pt}{4pt}{\sffamily\bfseries\large\color{ink}}}
\renewcommand{\subsection}{\Needspace{4\baselineskip}\@startsection{subsection}{2}{0pt}{9pt plus 1pt minus 1pt}{3pt}{\sffamily\bfseries\normalsize\color{ink}}}
\makeatother
\setcounter{secnumdepth}{0}
\begin{document}
{\sffamily\small\color{muted}Tensor computations\par}
\vspace{3pt}
{\raggedright\hyphenpenalty=10000\exhyphenpenalty=10000\sffamily\bfseries\fontsize{21}{24}\selectfont\color{ink}Near-optimal
type-2 bound for Gaussian sums of symmetric tensors\par}
\vspace{7pt}
{\sffamily\small\textbf{Difficulty:} extreme\quad\textbf{Importance:} broadly
interesting\par}
{\sffamily\small\bfseries\color{accent}Status: Partially resolved\par}
\vspace{4pt}
{\color{accent}\hrule height 0.8pt}
\vspace{6pt}
\textbf{Rating rationale:} Extreme because the low-p regime encounters
an explicit geometric barrier in tensor concentration; broad importance
spans probability, analysis, theoretical computer science and noisy
tensor computation.

\subsection{Statement}\label{statement}

For a symmetric real tensor \(T\in(\mathbb R^d)^{\otimes r}\), define \[
\|T\|_{\mathcal I_p}=\max_{\|x\|_p\le1}|\langle T,x^{\otimes r}\rangle|.
\] For every integer \(r\ge2\) and real \(2\le p<\infty\), do constants
\(C_{r,p}>0\) and \(a_{r,p}\ge0\) exist such that, for every \(d\ge2\),
\(N\ge1\), and deterministic symmetric tensors
\(T_1,\ldots,T_N\in(\mathbb R^d)^{\otimes r}\), \[
\mathbb E\left\|\sum_{i=1}^N g_iT_i\right\|_{\mathcal I_p}
\le C_{r,p}\,d^{1/2-1/p}[\log(2+dN)]^{a_{r,p}}
\left(\sum_{i=1}^N\|T_i\|_{\mathcal I_p}^{2}\right)^{1/2},
\] where \(g_i\) are independent real \(N(0,1)\) random variables?
Constants and logarithmic exponents must be independent of \(d,N\) and
the tensors. The explicit logarithmic form spells out the source's
\(\widetilde O_{r,p}\) notation.

\subsection{Relevance}\label{relevance}

This would control the injective norm of tensor perturbations with
general covariance, extending a central matrix concentration principle
to multilinear computation and noisy tensor models.

\subsection{References}\label{references}

\begin{enumerate}
\def\labelenumi{\arabic{enumi}.}
\tightlist
\item
  A. S. Bandeira, D. Dmitriev, K. Lucca, P. Nizić-Nikolac, and A.
  Rödder, \emph{Randomstrasse101: Open Problems of 2025},
  arXiv:2603.29571, 2026.
  \href{https://arxiv.org/html/2603.29571v1}{Primary text}, Entry 8 by
  K. Lucca, Conjecture 16 and equations (1)--(2).
\item
  A. S. Bandeira, S. Gopi, H. Jiang, K. Lucca, and T. Rothvoss, \emph{A
  Geometric Perspective on the Injective Norm of Sums of Random
  Tensors}. \href{https://arxiv.org/abs/2411.10633}{Primary preprint},
  introduction and tensor type-2 bounds; published as \emph{Tensor
  Concentration Inequalities: A Geometric Approach}, STOC 2025.
\end{enumerate}

\subsection{Status check --- 2026-09-10}\label{status-check-2026-09-10}

Rechecked \href{https://arxiv.org/html/2603.29571v1}{Randomstrasse101,
Entry 8, Conjecture 16} and the
\href{https://arxiv.org/abs/2411.10633}{underlying tensor-concentration
paper}, then searched for later type-2 bounds. The displayed inequality
is proved for p at least twice the tensor order and for the matrix case
r=p=2; the general low-p regime remains unresolved in the August 2026
manuscript. No full resolution was located. Restrictions on covariance
or summands do not establish the universal claim.
\end{document}
